\documentclass[12pt,a4paper]{article}
\usepackage[UTF8]{ctex}
\usepackage{amsmath,amssymb,amsthm}
\usepackage{geometry}

\geometry{margin=2.5cm}

\title{机械臂避免奇异的方法：SVD和det($J^T J$)}
\author{第5章和第11章补充说明}
\date{}

\begin{document}

\maketitle

\section{问题提出}

\textbf{问题}：机械臂避免奇异的方式有SVD和det($J^T J$)两个方式吗？分别介绍。

\textbf{简短回答}：是的，这两种方法都可以用于避免或处理奇异点，但它们的作用和实现方式不同。

\section{奇异点的定义和问题}

\subsection{奇异点的定义}

\textbf{运动学奇异点}：当雅可比矩阵$J(\theta)$的秩小于其最大可能秩时，机器人处于奇异配置。

\textbf{数学定义}：
\begin{itemize}
\item 对于6自由度机器人，$J(\theta) \in \mathbb{R}^{6 \times n}$
\item 如果$\text{rank}(J(\theta)) < \min(6, n)$，则机器人处于奇异点
\item 对于非冗余机器人（$n = 6$），如果$\det(J(\theta)) = 0$，则处于奇异点
\end{itemize}

\subsection{奇异点的问题}

\textbf{问题1：逆运动学不稳定}
\begin{itemize}
\item 在奇异点，$J(\theta)$不可逆
\item 无法直接计算$\dot{\theta} = J^{-1}(\theta)V$
\item 需要使用伪逆$J^{\dagger}(\theta)$
\end{itemize}

\textbf{问题2：关节速度过大}
\begin{itemize}
\item 在奇异点附近，即使末端执行器速度很小，关节速度也可能很大
\item 因为$J^{\dagger}(\theta)$的元素变得很大
\end{itemize}

\textbf{问题3：控制不稳定}
\begin{itemize}
\item 在奇异点附近，控制律可能不稳定
\item 任务空间质量矩阵$\Lambda(\theta) = J^{-T}MJ^{-1}$变得很大
\end{itemize}

\section{方法1：使用det($J^T J$)检测和避免奇异}

\subsection{基本原理}

\textbf{关键思想}：使用$\det(J^T J)$作为奇异性的度量。

\textbf{数学原理}：

对于非冗余机器人（$n = 6$），$J \in \mathbb{R}^{6 \times 6}$：
\begin{equation}
\det(J^T J) = \det(J)^2
\end{equation}

\textbf{性质}：
\begin{itemize}
\item 如果$\det(J^T J) = 0$，则$\det(J) = 0$，机器人处于奇异点
\item 如果$\det(J^T J)$很小，则接近奇异点
\item $\det(J^T J) \geq 0$（总是非负）
\end{itemize}

对于冗余机器人（$n > 6$），$J \in \mathbb{R}^{6 \times n}$：
\begin{itemize}
\item $J^T J \in \mathbb{R}^{n \times n}$
\item 如果$\det(J^T J) = 0$，则$J$不是满秩，机器人处于奇异点
\item $\det(J^T J)$仍然可以作为奇异性的度量
\end{itemize}

\subsection{实现方法}

\textbf{步骤1：计算$\det(J^T J)$}

\begin{equation}
d = \det(J^T J)
\end{equation}

\textbf{步骤2：设定阈值}

设定一个小的正数$\epsilon > 0$作为阈值：
\begin{equation}
\text{如果 } d < \epsilon \quad \Rightarrow \quad \text{接近奇异点}
\end{equation}

\textbf{步骤3：避免策略}

\textbf{策略1：阻尼伪逆（Damped Least Squares）}

当检测到接近奇异点时，使用阻尼伪逆：
\begin{equation}
J^{\dagger} = (J^T J + \lambda I)^{-1}J^T \tag{阻尼伪逆}
\end{equation}

其中$\lambda$是阻尼系数，可以根据$\det(J^T J)$自适应调整：
\begin{equation}
\lambda = \begin{cases}
0 & \text{如果 } d > d_{\text{threshold}} \\
\lambda_0\left(1 - \frac{d}{d_{\text{threshold}}}\right) & \text{如果 } d \leq d_{\text{threshold}}
\end{cases}
\end{equation}

其中$\lambda_0$是最大阻尼系数，$d_{\text{threshold}}$是阈值。

\textbf{策略2：调整期望速度}

当检测到接近奇异点时，减小期望速度：
\begin{equation}
V_d = \begin{cases}
V_{\text{desired}} & \text{如果 } d > d_{\text{threshold}} \\
V_{\text{desired}} \cdot \frac{d}{d_{\text{threshold}}} & \text{如果 } d \leq d_{\text{threshold}}
\end{cases}
\end{equation}

\textbf{策略3：改变轨迹}

当检测到接近奇异点时，改变运动轨迹，避开奇异点。

\subsection{优势和劣势}

\textbf{优势}：
\begin{itemize}
\item \textbf{计算简单}：只需要计算$\det(J^T J)$，计算复杂度$O(n^3)$
\item \textbf{直观}：$\det(J^T J)$直接反映奇异性
\item \textbf{易于实现}：可以实时计算和监控
\end{itemize}

\textbf{劣势}：
\begin{itemize}
\item \textbf{不能直接给出伪逆}：只能检测，不能直接计算稳定的伪逆
\item \textbf{需要额外计算}：检测到奇异后，还需要计算阻尼伪逆
\item \textbf{阈值选择}：需要选择合适的阈值$\epsilon$和$d_{\text{threshold}}$
\end{itemize}

\section{方法2：使用SVD（奇异值分解）避免奇异}

\subsection{基本原理}

\textbf{关键思想}：使用SVD分解雅可比矩阵，直接计算稳定的伪逆。

\textbf{SVD分解}：

对于雅可比矩阵$J \in \mathbb{R}^{m \times n}$（$m = 6$，$n \geq 6$），SVD分解为：
\begin{equation}
J = U\Sigma V^T
\end{equation}

其中：
\begin{itemize}
\item $U \in \mathbb{R}^{m \times m}$：左奇异向量矩阵（正交矩阵）
\item $\Sigma \in \mathbb{R}^{m \times n}$：奇异值矩阵（对角矩阵）
\begin{equation}
\Sigma = \begin{bmatrix}
\sigma_1 & 0 & \cdots & 0 \\
0 & \sigma_2 & \cdots & 0 \\
\vdots & \vdots & \ddots & \vdots \\
0 & 0 & \cdots & \sigma_r \\
0 & 0 & \cdots & 0 \\
\vdots & \vdots & \ddots & \vdots \\
0 & 0 & \cdots & 0
\end{bmatrix}
\end{equation}
其中$\sigma_1 \geq \sigma_2 \geq \cdots \geq \sigma_r > 0$是奇异值，$r = \text{rank}(J)$
\item $V \in \mathbb{R}^{n \times n}$：右奇异向量矩阵（正交矩阵）
\end{itemize}

\subsection{伪逆的计算}

\textbf{标准伪逆}：

使用SVD，伪逆可以直接计算为：
\begin{equation}
J^{\dagger} = V\Sigma^{\dagger}U^T
\end{equation}

其中$\Sigma^{\dagger}$是$\Sigma$的伪逆：
\begin{equation}
\Sigma^{\dagger} = \begin{bmatrix}
1/\sigma_1 & 0 & \cdots & 0 & 0 & \cdots & 0 \\
0 & 1/\sigma_2 & \cdots & 0 & 0 & \cdots & 0 \\
\vdots & \vdots & \ddots & \vdots & \vdots & \ddots & \vdots \\
0 & 0 & \cdots & 1/\sigma_r & 0 & \cdots & 0 \\
0 & 0 & \cdots & 0 & 0 & \cdots & 0 \\
\vdots & \vdots & \ddots & \vdots & \vdots & \ddots & \vdots \\
0 & 0 & \cdots & 0 & 0 & \cdots & 0
\end{bmatrix}
\end{equation}

\textbf{阻尼伪逆（Damped Least Squares）}：

当接近奇异点时（小的奇异值$\sigma_i$），使用阻尼伪逆：
\begin{equation}
\Sigma^{\dagger}_{\text{damped}} = \begin{bmatrix}
\frac{\sigma_1}{\sigma_1^2 + \lambda} & 0 & \cdots & 0 \\
0 & \frac{\sigma_2}{\sigma_2^2 + \lambda} & \cdots & 0 \\
\vdots & \vdots & \ddots & \vdots \\
0 & 0 & \cdots & \frac{\sigma_r}{\sigma_r^2 + \lambda} \\
0 & 0 & \cdots & 0 \\
\vdots & \vdots & \ddots & \vdots \\
0 & 0 & \cdots & 0
\end{bmatrix}
\end{equation}

阻尼伪逆：
\begin{equation}
J^{\dagger}_{\text{damped}} = V\Sigma^{\dagger}_{\text{damped}}U^T
\end{equation}

\subsection{自适应阻尼}

\textbf{自适应阻尼系数}：

根据最小奇异值$\sigma_{\min}$自适应调整阻尼系数：
\begin{equation}
\lambda = \begin{cases}
0 & \text{如果 } \sigma_{\min} > \sigma_{\text{threshold}} \\
\lambda_0\left(1 - \frac{\sigma_{\min}}{\sigma_{\text{threshold}}}\right)^2 & \text{如果 } \sigma_{\min} \leq \sigma_{\text{threshold}}
\end{cases}
\end{equation}

其中$\lambda_0$是最大阻尼系数，$\sigma_{\text{threshold}}$是阈值。

\textbf{物理意义}：
\begin{itemize}
\item 当$\sigma_{\min}$很大时，不需要阻尼（$\lambda = 0$）
\item 当$\sigma_{\min}$很小时，需要大的阻尼（$\lambda$接近$\lambda_0$）
\item 阻尼系数平滑变化，避免突然切换
\end{itemize}

\subsection{实现步骤}

\textbf{步骤1：计算SVD}
\begin{equation}
J = U\Sigma V^T
\end{equation}

\textbf{步骤2：检测奇异性}
\begin{equation}
\sigma_{\min} = \min\{\sigma_1, \sigma_2, \ldots, \sigma_r\}
\end{equation}

\textbf{步骤3：计算自适应阻尼}
\begin{equation}
\lambda = f(\sigma_{\min})
\end{equation}

\textbf{步骤4：计算阻尼伪逆}
\begin{equation}
J^{\dagger}_{\text{damped}} = V\Sigma^{\dagger}_{\text{damped}}U^T
\end{equation}

\textbf{步骤5：计算关节速度}
\begin{equation}
\dot{\theta} = J^{\dagger}_{\text{damped}}V
\end{equation}

\subsection{优势和劣势}

\textbf{优势}：
\begin{itemize}
\item \textbf{直接计算伪逆}：不需要先检测再计算，一步到位
\item \textbf{数值稳定}：SVD是数值稳定的算法
\item \textbf{自适应阻尼}：可以根据奇异值自动调整阻尼
\item \textbf{信息丰富}：SVD提供了奇异值和奇异向量，可以用于其他分析
\end{itemize}

\textbf{劣势}：
\begin{itemize}
\item \textbf{计算复杂}：SVD的计算复杂度是$O(\min(mn^2, m^2n))$，比直接求逆更复杂
\item \textbf{实时性}：对于高维矩阵，SVD可能较慢
\item \textbf{实现复杂}：需要实现SVD算法或使用库函数
\end{itemize}

\section{两种方法的对比}

\subsection{功能对比}

\begin{center}
\begin{tabular}{|l|l|l|}
\hline
\textbf{方面} & \textbf{det($J^T J$)} & \textbf{SVD} \\
\hline
主要功能 & 检测奇异性 & 计算稳定伪逆 \\
\hline
计算复杂度 & $O(n^3)$ & $O(\min(mn^2, m^2n))$ \\
\hline
数值稳定性 & 中等 & 高 \\
\hline
信息量 & 单一标量 & 完整分解 \\
\hline
实现难度 & 简单 & 中等 \\
\hline
\end{tabular}
\end{center}

\subsection{使用场景}

\textbf{det($J^T J$)方法适合}：
\begin{itemize}
\item 只需要检测奇异性
\item 计算资源有限
\item 需要快速判断
\item 可以接受额外的伪逆计算步骤
\end{itemize}

\textbf{SVD方法适合}：
\begin{itemize}
\item 需要直接计算稳定的伪逆
\item 需要完整的奇异值信息
\item 对数值稳定性要求高
\item 计算资源充足
\end{itemize}

\section{实际应用建议}

\subsection{混合方法（推荐）}

\textbf{策略}：结合两种方法的优势

\textbf{步骤1：快速检测（使用det($J^T J$)）}
\begin{equation}
d = \det(J^T J)
\end{equation}

\textbf{步骤2：判断是否需要SVD}
\begin{equation}
\text{如果 } d > d_{\text{threshold}} \quad \Rightarrow \quad \text{使用标准伪逆} \\
\text{如果 } d \leq d_{\text{threshold}} \quad \Rightarrow \quad \text{使用SVD计算阻尼伪逆}
\end{equation}

\textbf{步骤3：计算伪逆}
\begin{itemize}
\item 如果$d > d_{\text{threshold}}$：使用标准方法$J^{\dagger} = (J^T J)^{-1}J^T$
\item 如果$d \leq d_{\text{threshold}}$：使用SVD计算阻尼伪逆
\end{itemize}

\textbf{优势}：
\begin{itemize}
\item 大部分时间使用快速的标准方法
\item 只在接近奇异点时使用SVD
\item 平衡了计算效率和数值稳定性
\end{itemize}

\subsection{完整实现示例}

\textbf{伪代码}：

\begin{verbatim}
function compute_pseudoinverse(J, d_threshold, lambda_0):
    // 步骤1：快速检测
    d = det(J^T * J)
    
    // 步骤2：判断
    if d > d_threshold:
        // 远离奇异点，使用标准伪逆
        J_dagger = (J^T * J)^(-1) * J^T
    else:
        // 接近奇异点，使用SVD
        [U, Sigma, V] = svd(J)
        sigma_min = min(diag(Sigma))
        
        // 计算自适应阻尼
        if sigma_min > sigma_threshold:
            lambda = 0
        else:
            lambda = lambda_0 * (1 - sigma_min/sigma_threshold)^2
        
        // 计算阻尼伪逆
        Sigma_dagger_damped = diag(sigma_i / (sigma_i^2 + lambda))
        J_dagger = V * Sigma_dagger_damped * U^T
    
    return J_dagger
\end{verbatim}

\section{数学细节}

\subsection{det($J^T J$)与奇异值的关系}

\textbf{关系}：

对于$J \in \mathbb{R}^{m \times n}$的SVD分解$J = U\Sigma V^T$：
\begin{align}
J^T J &= (U\Sigma V^T)^T (U\Sigma V^T) \\
      &= V\Sigma^T U^T U\Sigma V^T \\
      &= V\Sigma^T\Sigma V^T \\
      &= V\Sigma^2 V^T
\end{align}

其中$\Sigma^2$是对角矩阵，对角线元素是$\sigma_1^2, \sigma_2^2, \ldots, \sigma_r^2$。

\textbf{行列式}：
\begin{equation}
\det(J^T J) = \det(V\Sigma^2 V^T) = \det(V)\det(\Sigma^2)\det(V^T) = \prod_{i=1}^r \sigma_i^2
\end{equation}

\textbf{结论}：
\begin{itemize}
\item $\det(J^T J) = \prod_{i=1}^r \sigma_i^2$
\item 如果$\det(J^T J) = 0$，则至少有一个奇异值为零
\item 如果$\det(J^T J)$很小，则最小奇异值很小
\end{itemize}

\subsection{条件数}

\textbf{条件数}：

矩阵$J$的条件数定义为：
\begin{equation}
\kappa(J) = \frac{\sigma_{\max}}{\sigma_{\min}} = \frac{\sigma_1}{\sigma_r}
\end{equation}

\textbf{物理意义}：
\begin{itemize}
\item 条件数大：接近奇异点
\item 条件数小：远离奇异点
\item 条件数$\geq 1$（总是）
\end{itemize}

\textbf{与det($J^T J$)的关系}：
\begin{equation}
\det(J^T J) = \prod_{i=1}^r \sigma_i^2 \leq \sigma_{\min}^{2r}
\end{equation}

因此：
\begin{equation}
\sigma_{\min} \geq \sqrt[r]{\det(J^T J)}
\end{equation}

\section{总结}

\subsection{关键要点}

\begin{enumerate}
\item \textbf{det($J^T J$)方法}：
\begin{itemize}
\item 主要用于检测奇异性
\item 计算简单，但需要额外的伪逆计算
\item 适合快速判断
\end{itemize}

\item \textbf{SVD方法}：
\begin{itemize}
\item 直接计算稳定的伪逆
\item 数值稳定，但计算复杂
\item 适合对稳定性要求高的场景
\end{itemize}

\item \textbf{推荐方案}：
\begin{itemize}
\item 使用det($J^T J$)快速检测
\item 只在接近奇异点时使用SVD
\item 结合两种方法的优势
\end{itemize}
\end{enumerate}

\subsection{公式总结}

\textbf{det($J^T J$)方法}：
\begin{equation}
d = \det(J^T J), \quad \lambda = f(d), \quad J^{\dagger} = (J^T J + \lambda I)^{-1}J^T
\end{equation}

\textbf{SVD方法}：
\begin{equation}
J = U\Sigma V^T, \quad \lambda = f(\sigma_{\min}), \quad J^{\dagger} = V\Sigma^{\dagger}_{\text{damped}}U^T
\end{equation}

\textbf{关系}：
\begin{equation}
\det(J^T J) = \prod_{i=1}^r \sigma_i^2
\end{equation}

\end{document}

